\documentclass[letters,10pt]{IEEEtran}
\usepackage{graphicx}
\usepackage{amsmath}
\usepackage{amssymb}
\usepackage{cite}
\usepackage{booktabs}
\usepackage{url}

\begin{document}

\title{Continuous-Time Neural Networks for Motor Imagery EEG Decoding: A Compact Study on Temporal Adaptivity and Interpretability}

\author{
  \IEEEauthorblockN{
    % Author(s) removed for blind review
  }
  \IEEEauthorblockA{
    % Institution removed for blind review
  }
}

\maketitle

\begin{abstract}
Motor imagery (MI) electroencephalography (EEG) signals exhibit non-stationary temporal dynamics that pose challenges for conventional discrete-time neural models. Continuous-time neural networks, particularly closed-form continuous-time (CfC) models, offer a promising alternative by directly modeling continuous-time dynamics through adaptive time constants. Despite the growing interest in continuous-time models, their effectiveness on multi-class MI-EEG decoding and the interpretability of their learned temporal representations remain underexplored. In this letter, we present a systematic evaluation of CfC against LSTM and EEGNet from a continuous-time modeling perspective on the BCI Competition IV Dataset 2a, and further test a lightweight Hybrid-CfC baseline that combines shallow EEGNet-style feature extraction with CfC dynamics. We implement CfC with per-unit, input-dependent time constants and revise the recurrent baselines with stronger temporal readout and a less aggressive early-stopping protocol. Experimental results show that pure CfC reaches 50.6\% accuracy and 0.493 macro F1, clearly outperforming LSTM at 41.5\% accuracy and 0.392 macro F1, while remaining slightly below EEGNet at 55.8\% accuracy and 0.540 macro F1. Hybrid-CfC obtains 54.1\% accuracy and 0.530 macro F1, and is not significantly different from EEGNet under paired subject-wise testing. Analysis of the learned time constant $\tau$ distributions reveals broader and mildly class-shifted temporal scales under the input-dependent parameterization, while substantial overlap remains across classes. Additionally, we evaluate model robustness under synthetic Gaussian noise perturbations. These findings indicate that continuous-time modeling can provide competitive MI-EEG decoding when the recurrent readout and training protocol are properly configured, and that explicit spatial inductive bias remains beneficial but not strictly necessary for effective performance.
\end{abstract}

\begin{IEEEkeywords}
Motor imagery EEG, closed-form continuous-time neural networks, brain-computer interface, temporal dynamics, interpretability
\end{IEEEkeywords}

\section{Introduction}
\addcontentsline{toc}{section}{Introduction}

Motor imagery (MI) based brain-computer interfaces (BCIs) translate covert movement intentions, as reflected in electroencephalography (EEG) signals, into actionable commands \cite{wolpaw2002brain}. Unlike evoked potentials, MI-EEG signals are self-paced and non-stationary, exhibiting complex temporal dynamics such as event-related desynchronization (ERD) and synchronization (ERS) in the $\mu$ (8--13 Hz) and $\beta$ (13--30 Hz) bands \cite{pfurtscheller2001motor}. These characteristics make MI-EEG decoding a challenging task that demands models capable of capturing fine-grained temporal dependencies.

In recent years, deep learning methods have become the dominant approach for MI-EEG classification. Convolutional neural networks (CNNs), particularly EEGNet \cite{lawhern2018eegnet}, and recurrent neural networks (RNNs), primarily long short-term memory (LSTM) networks, have demonstrated strong performance on BCI competition datasets \cite{schirrmeister2017deep}. However, these models operate on discrete time steps, which may not adequately capture the continuous-time dynamics inherent in EEG signals. LSTMs, for instance, process sequential data at fixed sampling intervals and rely on gating mechanisms rather than explicit continuous-time representations, potentially limiting their ability to adapt to the non-stationary characteristics of EEG signals. This mismatch motivates the exploration of continuous-time models for EEG decoding. At the same time, weak results from pure continuous-time recurrent modeling may indicate that temporal adaptivity alone is insufficient for MI-EEG, where spatial structure across electrodes also plays a central role.

Continuous-time neural networks offer a principled alternative by modeling the underlying dynamical system directly through differential equations. Liquid time-constant (LTC) networks \cite{hazir2019liquid} and their closed-form counterpart, closed-form continuous-time (CfC) networks \cite{oliveira2022closed}, parameterize the derivative of the hidden state as a function of the input, enabling adaptive time scales that can vary dynamically with the input. This property is particularly appealing for EEG signals, where the characteristic frequencies of neural oscillations differ across mental states and subjects.

Recent neuromorphic benchmark studies have begun to explore continuous-time models on motor imagery tasks. The Neuromorphic Sequential Arena benchmark \cite{chen2025neuromorphic} evaluated LTC networks on the OpenBMI dataset under a binary (left/right hand) MI setting, focusing primarily on efficiency metrics such as energy consumption, memory footprint, and training speed. While these results demonstrate the viability of continuous-time models for EEG-based MI decoding, several important questions remain unanswered. First, existing evaluations are limited to binary classification, whereas real-world BCI applications often require multi-class decoding (e.g., left hand, right hand, feet, and tongue in the standard four-class paradigm). Second, these benchmarks prioritize neuromorphic efficiency over classification performance and do not provide analysis of the learned temporal representations. Third, the robustness of continuous-time models to noise in EEG signals has not been systematically investigated.

To address these gaps, this letter presents a systematic study of CfC on multi-class MI-EEG decoding, with a focus on three aspects: classification performance, temporal interpretability, and noise robustness. Specifically, we make the following contributions:

\begin{itemize}
    \item We provide a systematic evaluation of CfC, LSTM, and EEGNet on the four-class BCI Competition IV Dataset 2a from a continuous-time modeling perspective, introducing time constant $\tau$ as an analysis tool.
    \item Through visualization and analysis of $\tau$ distributions across different MI classes, we examine the temporal scale learned by CfC and assess the extent to which it exhibits interpretable class-dependent behavior.
    \item We introduce a lightweight Hybrid-CfC baseline that combines shallow EEGNet-style spatial-temporal feature extraction with CfC dynamics to test whether continuous-time modeling benefits from explicit EEG spatial inductive bias.
    \item We evaluate the robustness of CfC under controlled Gaussian noise perturbations and analyze its performance degradation patterns compared to LSTM and EEGNet.
\end{itemize}

\section{Related Work}
\addcontentsline{toc}{section}{Related Work}

\subsection{MI-EEG Classification Methods}
Traditional MI-EEG classification relies on feature engineering methods, among which common spatial patterns (CSP) \cite{koles1990spatial} and its variants, such as filter bank CSP (FBCSP) \cite{ang2008filter}, have been the most widely used. These methods extract spatial filters that maximize the variance ratio between classes, followed by classifiers such as linear discriminant analysis (LDA). Despite their popularity, CSP-based methods are sensitive to non-stationarities in EEG signals and require careful subject-specific parameter tuning.

Deep learning approaches have largely superseded traditional methods in recent years. CNNs, particularly EEGNet \cite{lawhern2018eegnet}, apply spatial and temporal convolutions directly to raw EEG signals and have achieved competitive results across multiple BCI paradigms. RNNs, especially LSTM networks, model sequential dependencies through gated memory cells and have been widely adopted for MI-EEG decoding \cite{schirrmeister2017deep}. Transformer-based architectures have also been explored for MI-EEG \cite{song2022transformer}, though their high computational cost and data requirements limit practical deployment.

\subsection{Continuous-Time Neural Networks}
Continuous-time neural networks parameterize the dynamics of hidden states through ordinary differential equations (ODEs), allowing the model to adapt its time scale continuously rather than at fixed discrete steps. The liquid time-constant (LTC) network \cite{hazir2019liquid} defines the derivative of the hidden state as a nonlinear function of both the input and the current state, with a learnable time constant $\tau$ controlling the update rate. This formulation enables the network to exhibit adaptive behavior that varies with input characteristics.

The closed-form continuous-time (CfC) network \cite{oliveira2022closed} approximates the solution of the LTC ODE in closed form, achieving inference speeds one to five orders of magnitude faster than discretized ODE solvers while maintaining the core property of continuous-time dynamics. CfC has shown promising results in domains such as healthcare monitoring \cite{bhattacharjee2026adaptive} and speech processing \cite{oliveira2022closed}, but its application to motor imagery EEG remains underexplored.

\subsection{Neuromorphic Benchmark Work}
Recent efforts have evaluated continuous-time models within neuromorphic benchmarking frameworks. Chen et al. \cite{chen2025neuromorphic} introduced the Neuromorphic Sequential Arena (NSA), a comprehensive benchmark for neuromorphic temporal processing that includes LTC networks on EEG-MI tasks from the OpenBMI dataset. These evaluations, however, are primarily efficiency-oriented (energy, memory, and training speed) and employ a binary MI classification setting. In contrast to prior efficiency-oriented binary evaluations, this work investigates the more challenging four-class MI-EEG decoding task, with a focus on classification performance and temporal interpretability through $\tau$ analysis.

\section{Proposed Method}
\addcontentsline{toc}{section}{Proposed Method}

\subsection{Closed-Form Continuous-Time Model}
The core of CfC \cite{oliveira2022closed} is the closed-form approximation of the LTC ODE. In our implementation, the raw EEG sample at time step $t$ is first projected to $\mathbf{z}_t \in \mathbb{R}^{H}$, and the CfC cell updates the hidden state $\mathbf{h}_t \in \mathbb{R}^{H}$ through
\begin{align}
\tilde{\mathbf{h}}_t &= \tanh\!\left(\mathbf{W}_c[\mathbf{z}_t; \mathbf{h}_{t-1}] + \mathbf{b}_c\right), \\
\boldsymbol{\tau}_t &= \mathrm{softplus}\!\left(\mathbf{W}_{\tau,2}\,\phi\!\left(\mathbf{W}_{\tau,1}[\mathbf{z}_t; \mathbf{h}_{t-1}] + \mathbf{b}_{\tau,1}\right) + \mathbf{b}_{\tau,2}\right) + \epsilon, \\
\mathbf{h}_t &= \exp(-\Delta t / \boldsymbol{\tau}_t) \odot \mathbf{h}_{t-1}
+ \left(1 - \exp(-\Delta t / \boldsymbol{\tau}_t)\right) \odot \tilde{\mathbf{h}}_t ,
\end{align}
where $\phi(\cdot)$ is a SiLU activation, $\epsilon = 10^{-3}$ enforces positivity, and $\boldsymbol{\tau}_t \in \mathbb{R}^{H}$ is a per-unit, input-dependent time-constant vector. We additionally apply layer normalization to the concatenated $[\mathbf{z}_t; \mathbf{h}_{t-1}]$ before the candidate and $\tau$ heads. Compared with a single global learnable $\tau$, this parameterization allows the effective temporal scale to vary dynamically with both the current input and the current recurrent state, which is more faithful to the adaptive continuous-time formulation and more suitable for downstream $\tau$ analysis.

\subsection{Baseline Models}
As reference baselines, we adopt two established models. The LSTM \cite{hochreiter1997long} is the standard discrete-time recurrent baseline, using input, forget, and output gates to manage memory flow. The EEGNet \cite{lawhern2018eegnet} is a compact CNN designed specifically for EEG, using depthwise and separable convolutions to capture spatial and temporal features with few parameters. EEGNet is included as a representative compact CNN specifically designed for EEG decoding, serving as a strong spatial-temporal baseline. All models share the same preprocessing pipeline, optimizer family, and validation protocol; model sizes are reported explicitly rather than hidden behind unmatched capacity budgets.

\subsection{Hybrid Spatial-Continuous Baseline}
To examine whether the main weakness of pure CfC arises from missing spatial inductive bias rather than from continuous-time modeling itself, we introduce a lightweight Hybrid-CfC baseline. The hybrid model first applies a shallow EEGNet-style frontend with 8 temporal filters (kernel size $1 \times 64$), a depthwise spatial convolution spanning all 22 electrodes to produce 16 feature maps, and a compact separable temporal block composed of a depthwise $1 \times 16$ convolution and a pointwise $1 \times 1$ projection. Average pooling with kernel size $1 \times 4$ is used after each convolutional block, with dropout rate 0.25 in the frontend. The resulting feature sequence is then projected to a 128-dimensional CfC module and classified by a linear head. This design is intentionally minimal: it is not proposed as a new full architecture, but rather as a diagnostic baseline for testing whether explicit spatial feature extraction complements continuous-time dynamics in MI-EEG.

\subsection{Classification Pipeline}
The processing pipeline consists of four stages. Raw EEG signals are first band-pass filtered (8--30 Hz) to retain motor-related frequency bands. Filtered signals are then segmented into trials aligned to MI cues. Each model processes the preprocessed data to produce class predictions. Cross-entropy loss is used for training, and classification accuracy and macro-averaged F1 score are reported as evaluation metrics.

\begin{figure}[t]
  \centering
  \IfFileExists{pipeline_placeholder.pdf}{%
    \includegraphics[width=0.9\linewidth]{pipeline_placeholder.pdf}%
  }{%
    \fbox{\parbox[c][0.22\textheight][c]{0.86\linewidth}{\centering Pipeline figure placeholder}}%
  }
  \caption{Block diagram of the MI-EEG classification pipeline. EEG signals are band-pass filtered, segmented into trials, fed into the models (CfC, Hybrid-CfC, LSTM, or EEGNet), and classified into one of four motor imagery classes.}
  \label{fig:pipeline}
\end{figure}

\section{Experiments}
\addcontentsline{toc}{section}{Experiments}

\subsection{Dataset}
We evaluate on the BCI Competition IV Dataset 2a \cite{tangermann2012review}, a standard benchmark for four-class MI-EEG decoding. The dataset contains recordings from 9 subjects performing four MI tasks: left hand, right hand, feet, and tongue. Each subject performed 72 trials per class, recorded from 22 Ag/AgCl electrodes at 250 Hz. We use the standard motor-imagery window from 2 s to 6 s after cue onset. The evaluation follows a subject-dependent protocol in which the trials from both sessions are pooled and then split by 5-fold stratified cross-validation within each subject. We report the mean and standard deviation of accuracy and F1 score across all subjects.

\subsection{Experimental Setup}
All models are implemented in PyTorch. EEG signals are band-pass filtered between 8 Hz and 30 Hz using a Butterworth filter (order 4). After filtering, trials are downsampled by a factor of 2 to 125 Hz to reduce the cost of repeated cross-validation. No additional artifact rejection or manual curation is applied. For CfC, the hidden dimension is set to 128, and the time constant is parameterized per hidden unit by the input-dependent MLP in (2), with the final bias initialized to yield $\tau \approx 1.0$ at the start of training. Instead of using only the final recurrent state, both CfC and LSTM employ temporal mean-max pooling over the full hidden sequence before classification, which provides a stronger but still standard sequence-level readout. For LSTM, we use a two-layer network with 128 hidden units per layer and dropout 0.4 between recurrent layers and before classification. EEGNet follows the compact configuration of \cite{lawhern2018eegnet}. Hybrid-CfC uses the shallow EEGNet-style frontend described above followed by a CfC module with hidden size 128. The trainable parameter counts are 2.4k for EEGNet, 210.9k for LSTM, 86.8k for CfC, and 87.0k for Hybrid-CfC. In each fold, train/validation/test trials are strictly disjoint, and all model selection decisions are made using only the training split and its validation subset, reducing the risk of cross-fold leakage. All models are trained with the Adam optimizer, initial learning rate of $1\times10^{-3}$, weight decay of $1\times10^{-4}$, cosine annealing, and a batch size of 64 for up to 80 epochs. Early stopping uses validation accuracy with patience of 20 epochs after a 25-epoch warm-up period. We intentionally keep one shared preprocessing and optimization protocol across models and report parameter counts explicitly to improve fairness and reproducibility.

\subsection{Classification Results}
\begin{table}[t]
  \centering
  \caption{Classification accuracy ($\%$), macro F1 score, and trainable parameter counts on BCI Competition IV Dataset 2a (mean $\pm$ std across 9 subjects, 5-fold CV).}
  \begin{tabular}{lccc}
    \toprule
    \textbf{Model} & \textbf{Params} & \textbf{Accuracy (\%)} & \textbf{F1 Score} \\
    \midrule
    EEGNet \cite{lawhern2018eegnet} & 2.4k & 55.8 $\pm$ 18.2 & 0.540 $\pm$ 0.182 \\
    Hybrid-CfC & 87.0k & 54.1 $\pm$ 19.1 & 0.530 $\pm$ 0.198 \\
    CfC \cite{oliveira2022closed} & 86.8k & 50.6 $\pm$ 17.7 & 0.493 $\pm$ 0.185 \\
    LSTM \cite{hochreiter1997long} & 210.9k & 41.5 $\pm$ 12.6 & 0.392 $\pm$ 0.140 \\
    \bottomrule
  \end{tabular}
  \label{tab:results}
\end{table}

The results in Table \ref{tab:results} show that EEGNet remains the strongest model on average, but the gap to the recurrent baselines is now much smaller than in the earlier simplified setting. Pure CfC reaches 50.6\% accuracy and 0.493 macro F1, clearly above the 41.5\% accuracy and 0.392 macro F1 obtained by LSTM. Hybrid-CfC further improves to 54.1\% accuracy and 0.530 macro F1, approaching EEGNet while remaining within the same general parameter scale as pure CfC. These results indicate that the earlier near-chance recurrent results were not due to an inherently broken pipeline; rather, recurrent MI-EEG decoding is highly sensitive to sequence readout and training dynamics. Under a more appropriate recurrent setup, pure continuous-time modeling already becomes competitive on this task.

The large standard deviations in Table \ref{tab:results} reflect substantial subject-to-subject variability, which is typical for MI-EEG but still requires explicit reporting. To make this variance transparent, the supplementary material includes a subject-wise accuracy boxplot and a per-subject accuracy table. Importantly, CfC outperforms LSTM on all 9 subjects, indicating that the recurrent gain is not driven by a small subset of outliers. By contrast, Hybrid-CfC exceeds EEGNet on 5 of 9 subjects, which is consistent with the non-significant subject-wise difference between these two models.

We further conduct paired $t$-tests on subject-wise mean accuracies. CfC significantly outperforms LSTM (CfC vs. LSTM: $p = 0.003$, Cohen's $d = 1.373$), indicating that the continuous-time recurrent formulation provides a measurable advantage over a comparably sized discrete-time recurrent baseline under the same protocol. EEGNet still outperforms CfC by a modest but statistically significant margin (CfC vs. EEGNet: $p = 0.047$, Cohen's $d = -0.784$). However, neither the comparison between Hybrid-CfC and CfC (Hybrid-CfC vs. CfC: $p = 0.157$, Cohen's $d = 0.521$) nor the comparison between Hybrid-CfC and EEGNet (Hybrid-CfC vs. EEGNet: $p = 0.292$, Cohen's $d = -0.376$) reaches statistical significance. These results support a stronger conclusion than before: continuous-time modeling is effective in its own right on MI-EEG, while explicit spatial feature extraction remains helpful and can close most of the remaining gap to EEGNet.

\subsection{Interpretability Analysis: Time Constant Distributions}
A distinctive advantage of CfC over LSTM is the interpretability of its learned time constants. We extract the $\tau$ values from the CfC hidden states at each time step and analyze their distributions across the four MI classes.

\begin{figure}[t]
  \centering
  \IfFileExists{outputs/recurrent_fixed/tau_dist_placeholder.pdf}{%
    \includegraphics[width=0.9\linewidth]{outputs/recurrent_fixed/tau_dist_placeholder.pdf}%
  }{%
    \fbox{\parbox[c][0.22\textheight][c]{0.86\linewidth}{\centering Time-constant distribution placeholder}}%
  }
  \caption{Distribution of learned time constant $\tau$ across four MI classes. Each subplot shows the histogram of $\tau$ values for one class (left hand, right hand, feet, tongue). The color indicates the density of time constants observed at each value.}
  \label{fig:tau}
\end{figure}

Figure \ref{fig:tau} shows that the class-wise $\tau$ distributions remain overlapping, but now exhibit broader dispersion and modest class-wise shifts under the input-dependent parameterization. The median $\tau$ values differ across classes (left hand: 1.412, right hand: 1.326, feet: 1.456, tongue: 1.455), with right hand tending toward slightly shorter temporal constants than feet and tongue. Because $\tau$ is generated dynamically from the current input and hidden state, these histograms reflect actual sample-wise temporal adaptation rather than a single static learnable constant. The substantial overlap between the distributions indicates that $\tau$ should still be interpreted as a soft temporal descriptor rather than a strongly class-separating biomarker, but the broader range suggests that the revised CfC does make more active use of adaptive temporal scaling than the earlier simplified implementation.

\subsection{Noise Robustness}
To evaluate robustness, we superimpose Gaussian noise onto the preprocessed EEG signals at three signal-to-noise ratio (SNR) levels: 10 dB, 5 dB, and 0 dB. Table \ref{tab:noise} reports the mean accuracy across all subjects under each condition.

\begin{table}[t]
  \centering
  \caption{Classification accuracy ($\%$) under Gaussian noise with varying SNR levels (mean $\pm$ std across 9 subjects).}
  \scriptsize
  \begin{tabular}{lccc}
    \toprule
    \textbf{Model} & \textbf{SNR = 10 dB} & \textbf{SNR = 5 dB} & \textbf{SNR = 0 dB} \\
    \midrule
    EEGNet & 47.3 $\pm$ 18.0 & 33.8 $\pm$ 7.7 & 25.9 $\pm$ 1.1 \\
    Hybrid-CfC & 45.6 $\pm$ 16.4 & 32.4 $\pm$ 6.8 & 25.8 $\pm$ 1.2 \\
    CfC & 47.7 $\pm$ 16.7 & 40.6 $\pm$ 12.1 & 27.5 $\pm$ 3.2 \\
    LSTM & 41.6 $\pm$ 12.5 & 41.1 $\pm$ 11.7 & 39.6 $\pm$ 10.4 \\
    \bottomrule
  \end{tabular}
  \label{tab:noise}
\end{table}

The results indicate different robustness profiles across the models. EEGNet and Hybrid-CfC retain strong performance at 10 dB SNR but degrade substantially as the noise level increases. CfC is particularly competitive under moderate noise, slightly exceeding EEGNet at 10 dB and maintaining 40.6\% accuracy at 5 dB. To avoid over-interpreting these mean values, we additionally summarize subject-wise variability and uncorrected paired tests. The CfC-EEGNet difference is not significant at 10 dB ($p = 0.790$) or 0 dB ($p = 0.191$), but CfC is significantly higher at 5 dB ($p = 0.022$). Hybrid-CfC does not differ significantly from EEGNet at any tested SNR (all $p > 0.19$). LSTM exhibits the flattest degradation curve and the highest mean accuracy at 0 dB, although the corresponding standard deviations remain large. Overall, the recurrent models appear less sensitive to synthetic Gaussian perturbations than the CNN-heavy baselines, but the noise analysis should be interpreted as robustness evidence under a synthetic perturbation model rather than as a definitive claim about real-world artifact tolerance.

\section{Discussion and Conclusion}
\addcontentsline{toc}{section}{Discussion and Conclusion}

This letter presents a systematic evaluation of closed-form continuous-time neural networks for motor imagery EEG decoding. Our experiments on the four-class BCI Competition IV Dataset 2a show that, with per-unit input-dependent time constants, stronger sequence readout, and a less aggressive early-stopping policy, pure CfC becomes a genuinely effective baseline rather than a near-chance one. In the final protocol, CfC significantly outperforms LSTM and reaches performance reasonably close to EEGNet, while Hybrid-CfC further narrows the gap and is not significantly different from EEGNet. These results substantially strengthen the central claim of the paper: continuous-time modeling is useful for MI-EEG on its own, and can be further complemented by explicit spatial inductive bias.

The revised $\tau$ analysis also paints a more credible interpretability picture. While the class-wise $\tau$ distributions are still overlapping, they now exhibit broader dispersion and mild class-dependent shifts, indicating that the model is making nontrivial use of adaptive temporal scaling rather than collapsing to an almost constant time scale. Taken together, the evidence suggests that the main limitation of continuous-time modeling in MI-EEG is not ineffectiveness per se, but sensitivity to architectural and training choices that determine whether temporal adaptivity is actually expressed.

\textbf{Limitations.} This study has several limitations. First, we evaluate on a single dataset; generalization across datasets and subjects requires further investigation. Second, the number of baseline models is limited; incorporating a broader range of methods would strengthen the comparison. Third, although we now report subject-wise dispersion and supplementary per-subject results, the variance across subjects remains substantial. Fourth, the noise perturbation experiments use synthetic Gaussian noise, which may not fully replicate real-world EEG artifacts such as eye blinks, muscle activity, or electrode motion. Fifth, the interpretability analysis is based on learned $\tau$ distributions and does not make direct neuroscientific claims about neural mechanisms.

\textbf{Future Work.} Future directions include (i) leave-one-subject-out evaluation to assess cross-subject generalization, (ii) validation on real-world EEG artifacts from clinical or ambulatory settings, and (iii) systematic exploration of hybrid spatial-continuous architectures for real-time BCI applications where the low computational cost of CfC may provide a practical advantage.

\bibliographystyle{IEEEtran}
\bibliography{references}

\end{document}
